\documentclass[11pt,a4paper]{article}
\usepackage[margin=27mm]{geometry}
\usepackage{fontspec,unicode-math}
\setmainfont{STIXGeneral}
\setmathfont{STIX Two Math}
\usepackage{amsmath,amsthm,booktabs,longtable,hyperref}
\setlength{\emergencystretch}{3em}
\newtheorem{theorem}{Theorem}
\newtheorem{lemma}[theorem]{Lemma}
\newtheorem{proposition}[theorem]{Proposition}
\newtheorem{corollary}[theorem]{Corollary}
\theoremstyle{remark}\newtheorem{remark}[theorem]{Remark}
\newcommand{\Aut}{\operatorname{Aut}}
\newcommand{\id}{\operatorname{id}}
\newcommand{\NN}{\mathbb N}
\newcommand{\ZZ}{\mathbb Z}
\newcommand{\Sym}{\operatorname{Sym}}
\newcommand{\W}{\mathsf W}
\hypersetup{colorlinks=true,linkcolor=blue,urlcolor=blue}
\newcommand{\gen}[1]{\langle #1\rangle}

\title{Coset intersections under group operations}
\author{Human Compression Project}\date{}
\begin{document}\maketitle
\begin{abstract}
Common representatives of incident cosets provide a uniform language for free products, coherent amalgams, HNN extensions and orbit constructions. We prove separate transfer theorems for tree-supported cosets and semidirect-product fibers, retaining nontrivial chamber stabilizers. Exact panel and residue formulas distinguish the hypotheses needed for firmness, thinness and residual connectedness. Cyclic markings and inverse presentation maps give the Artin applications, while orbital presentations and module changes of coordinates retain the complete shared theory used by the diagonal and wreath constructions. The accompanying application applies these results to Piedade and Tranchida's 2025 paper with explicit source qualifications.
\end{abstract}
The common thread is the existence of a representative for every finite family of pairwise intersecting parabolic cosets. On a tree, representatives are transported through full subcomplexes or compatible local cosets; in a semidirect product they are assembled in a common second-coordinate fiber. These mechanisms have different hypotheses. No theorem about arbitrary graphs of geometries follows simply by combining them.

This is a complete written theory paper for three existing application contracts: N005 (amalgamation, HNN extensions and twisting), N006 (wreath and semidirect constructions) and N007 (element-diagonal constructions). Its modular sections have explicit uses. The signed-module section is used by N006, not by N005; the tree and cyclic sections supply the additional N005 routes. Full relevant common theory remains part of any eventual comparison. This draft contains no length comparison or new theorem-credit claim.
\section{Cosets, flags and residue generation}\label{O-sec:cosets}
All type sets in this paper are finite. Cosets are left cosets. In a coset geometry $(B,(B_i)_{i\in I})$ write $B_J=\bigcap_{i\in J}B_i$, with $B_\varnothing=B$. A regular hypertope is thin, residually connected and flag-transitive. Automorphisms preserve types. In its faithful standard model, the group acts freely and transitively on chambers. Its distinguished involutions $b_i$ interchange a fixed chamber and its $i$-adjacent chamber; then
\begin{equation}\label{O-eq:parabolic}
B_J=\langle b_i:i\in I\setminus J\rangle,\qquad B_I=1.
\end{equation}
We recall why the geometric hypotheses give these assertions. In a thin geometry a type-preserving automorphism fixing a chamber fixes all its adjacent chambers. Residual connectivity implies that the chamber graph, and the chamber graph of every residue, is connected: by induction on rank, chambers containing a given element are connected; a path in the incidence graph connects these families, using chambers containing each consecutive incident pair. The rank-one and rank-two cases start the induction. Thus a chamber stabilizer is trivial. The same argument inside the residue of the base flag of types $J$ proves (\ref{O-eq:parabolic}). It also proves that the full type-preserving automorphism group is the displayed group whenever a faithful group already acts transitively on chambers.

For a second regular hypertope $(A,(A_x)_{x\in S})$, put $A(T)=\langle a_x:x\in T\rangle$ for $T\subseteq S$. Formula (\ref{O-eq:parabolic}) gives
\begin{equation}\label{O-eq:ip}
A(T)\cap A(T')=A(T\cap T').
\end{equation}
We will use both this intersection property and flag-transitivity. They are different requirements: the intersection property alone is not a general criterion for a nonstring coset system to be a geometry.

\begin{lemma}\label{F-lem:cosets}
A coset incidence system $(A,(A_i)_{i\in I})$ is flag-transitive if
and only if the family of all left cosets of all $A_J$, $J\subseteq I$,
is Helly.
\end{lemma}
\begin{proof}
A flag is a collection of pairwise intersecting cosets of distinct
maximal parabolics. It is a translate of a standard flag precisely
when these cosets have a common element. This proves one implication.
For the other, expand each coset $aA_J$ into the intersection of
$aA_j$, $j\in J$. If two original cosets intersect, every maximal
coset from one meets every maximal coset from the other. Within one
expanded coset they meet at $a$. Two such cosets of the same maximal
parabolic are therefore equal. The resulting distinct-type collection
is a flag, and flag transitivity supplies a common element. Empty
$J$ imposes no constraint. A standard flag extends to the standard
chamber, so flag transitivity also supplies the geometric property.
\end{proof}

\begin{lemma}[Coset intersection]\label{O-lem:helly}
If $B$ is flag-transitive on its coset geometry, every finite pairwise-intersecting family of cosets of standard parabolic subgroups has nonempty total intersection. In particular, for $H,K,J\subseteq I$,
\[
B_HB_J\cap B_KB_J=B_{H\cup K}B_J.
\]
\end{lemma}
\begin{proof}
Lemma~\ref{F-lem:cosets} gives the first assertion. For the last assertion, if $g\in B_HB_J\cap B_KB_J$, the three cosets $B_H,B_K,gB_J$ intersect pairwise. A point of their common intersection belongs to $B_{H\cup K}$ and to $gB_J$. The reverse inclusion follows from subgroup containment.
\end{proof}

\begin{theorem}[Element orbits]\label{O-thm:elements}
Fix $j\in I$ and $x_0=B_j\in X=B/B_j$. If $B$ is flag-transitive, then
\begin{equation}\label{O-eq:orbit}
(B_Hx_0)\cap(B_Kx_0)=B_{H\cup K}x_0\qquad(H,K\subseteq I).
\end{equation}
Thus the complementary notation $O(M)=B_{I\setminus M}x_0$ satisfies $O(M)\cap O(N)=O(M\cap N)$.
\end{theorem}
\begin{proof}
The inverse image of $B_Hx_0$ under $B\to B/B_j$ is $B_HB_j$. Apply Lemma~\ref{O-lem:helly} and take images. Finally $(I\setminus M)\cup(I\setminus N)=I\setminus(M\cap N)$.
\end{proof}





\subsection{Parabolics with a nontrivial chamber stabilizer}
We now allow flag-transitive coset geometries whose acting groups are not
faithful. The type sets remain finite and nonempty. In a geometry
$\alpha=(A,(A_x)_{x\in S})$ put $A_E=\bigcap_{x\in E}A_x$, with
$A_\varnothing=A$, and call $A_S$ the Borel subgroup. It need not be
trivial. The residue of the standard flag of types $E$ is the coset
geometry of $A_E$ with parabolics $A_{E\cup\{x\}}$, $x\notin E$.
To see this identification, take an element incident with that flag;
flag transitivity supplies a common representative in $A_E$. Two such
representatives define the same element exactly when their quotient
lies in $A_E\cap A_x$. This also identifies incidence by intersecting
the corresponding cosets.

For a coset geometry of rank at least two the incidence graph is
connected if and only if its maximal parabolics generate the acting
group. Indeed, if $H$ is that generated subgroup, paths from the
standard flag remain among cosets with representatives in $H$.
Conversely multiplication successively by elements of the parabolics
provides paths through incident cosets; when two successive steps have
the same type, an intervening different type gives the same conclusion.
Thus the components are the translates of this component. Applied to
every residue, this proves that residual connectedness is equivalent to
\begin{equation}\label{O-eq:rc-borel}
 A_E=\langle A_{E\cup\{x\}}:x\notin E\rangle
 \quad\text{if }|S\setminus E|\geq2.
\end{equation}
For rank one residual connectedness imposes no condition.

\begin{lemma}[Generation of parabolics]\label{O-lem:parabolic-join}
If $\alpha$ is residually connected, then
\[
 A_E=\langle A_S,\ A_{S\setminus\{x\}}:x\notin E\rangle.
\]
Consequently, for any nonempty family $(E_\nu)$ of subsets of $S$,
\begin{equation}\label{O-eq:parabolic-join}
 \langle A_{E_\nu}:\nu\rangle=A_{\bigcap_\nu E_\nu}.
\end{equation}
\end{lemma}
\begin{proof}
Induct on $|S\setminus E|$. For zero this is $A_S=A_S$; for one
the displayed generator is $A_E$ itself. For larger complements,
use (\ref{O-eq:rc-borel}) and the induction hypothesis on each child.
Every missing type remains missing in at least one child, so all and
only the indicated minimal parabolics occur. For the second assertion,
each group on the left lies in the right-hand group. Conversely, every
minimal parabolic required to generate the right-hand group occurs
in some $A_{E_\nu}$, and $A_S$ lies in every one. The nonempty-family
hypothesis is necessary when $A_S\ne1$.
\end{proof}
The explicit $A_S$ in the first formula must not be dropped at
$E=S$. This avoids the empty-generated-subgroup problem in the
source's background equivalences.

\section{Tree gluing and coherent amalgamation}\label{sec:trees}


An abstract simplicial complex is \emph{flag} if each finite set of pairwise
adjacent vertices spans a simplex. A subcomplex is \emph{full} if it contains
every simplex of the ambient complex whose vertices it contains. We first
isolate a standard elementary gluing principle and prove the form needed
here. No novelty is claimed for the principle itself.

\begin{lemma}\label{T-lem:tree}
Let $T$ be a tree. For each vertex $v$ let $X_v$ be a flag simplicial
complex. For each edge $e=uv$ let $X_e$ be a complex equipped with
injective simplicial maps to $X_u$ and $X_v$, both with full image.
Form $X$ by identifying these images. Then each $X_v$ embeds as a full
subcomplex of $X$, and $X$ is flag.
\end{lemma}
\begin{proof}
An identification of vertices is a finite sequence of identifications
across edges of $T$. Whenever the underlying walk immediately crosses an
edge and returns, the two steps cancel, because the edge maps are
injective. Reducing the walk proves that a vertex of one $X_v$ cannot
be identified with a different vertex of the same $X_v$. It also proves
that the set $S(x)$ of vertices $v$ for which $x$ has a representative
in $X_v$ is a connected subtree: representatives along the unique path
between any two such vertices are forced and exist.

If two global vertices $x,y$ are represented in both $X_u$ and $X_v$,
they are represented in every block along the path from $u$ to $v$.
If they are adjacent in one of these blocks, fullness of the first edge
image puts their edge in that edge complex, and hence in the next block.
Repeating this argument transports adjacency along the path. More
generally, the identical argument transports every finite simplex.
Consequently every block is full in the union.

Now let $x_1,\ldots,x_m$ be pairwise adjacent vertices of $X$.
Their support subtrees $S(x_i)$ meet pairwise, since every edge of $X$
comes from a block. Finite families of pairwise intersecting subtrees
of a tree have a common vertex. Here is an explicit proof of that fact.
Root $T$ at any vertex. In each nonempty subtree take the unique vertex
nearest the root, and choose one, $p$, of greatest distance among these
finitely many nearest vertices. For another subtree choose a point where
it meets the subtree associated to $p$. The two nearest vertices lie on
the rooted path to that point; the deeper one is $p$. Connectedness of
the other subtree therefore puts $p$ in it. Thus $p$ lies in every subtree.

All $x_i$ are consequently represented in $X_p$. Adjacency transports
to that block as above. Since $X_p$ is flag, these representatives span
a simplex in $X_p$ and therefore in $X$.
\end{proof}

If $Y$ is flag and $\sigma$ is a simplex, its closed star is a full
subcomplex: its vertices are those of $\sigma$ and all vertices adjacent
to every vertex of $\sigma$. Any simplex on these vertices joins
$\sigma$ to a clique, and the flag property fills that clique.
This observation supplies the full edge inclusions in the next result.

\begin{theorem}\label{T-thm:amalgam}
Let $\alpha=(A,(A_i)_{i\in I_A})$ and
$\beta=(B,(B_i)_{i\in I_B})$ be flag-transitive coset geometries
of finite rank. Write $L=I_A\cap I_B$, $P=I_A\setminus L$ and
$Q=I_B\setminus L$. Suppose a fixed isomorphism identifies
$C=A_P$ with $B_Q$ and identifies $A_{P\cup J}$ with $B_{Q\cup J}$
for every $J\subseteq L$.
In $G=A*_C B$ put
\[
 G_i=\langle A_{\{i\}\cap I_A},B_{\{i\}\cap I_B}\rangle
 \quad(i\in I_A\cup I_B).
\]
Then the coset incidence system $(G,(G_i))$ is a flag-transitive
geometry. Its chamber stabilizer is the common subgroup
$H=A_{I_A}=B_{I_B}$, and for every $J\subseteq I_A\cup I_B$,
\[
 G_J=A_{J\cap I_A}*_{C\cap A_{J\cap I_A}}B_{J\cap I_B}.
\]
\end{theorem}
\begin{proof}
We use the usual tree of the amalgam: vertices are left cosets $gA,gB$,
edges are $gC$, and $gC$ joins $gA$ to $gB$. The reduced word theorem
proves that this graph is a tree. Using left cosets here only changes
the convention by inversion from the usual right-coset geometry.

At an $A$ vertex put the flag complex of $\alpha$ joined to a simplex
whose types are $Q$. At a $B$ vertex put the flag complex of $\beta$
joined to a simplex of types $P$. At an edge put the flag complex of
the common residue
\[
 (C,(C\cap A_l)_{l\in L})=(C,(C\cap B_l)_{l\in L})
\]
joined to the simplex of types $P\cup Q$. Flag transitivity identifies
this residue with the link of the standard $P$ face in $\alpha$ and
with the link of the standard $Q$ face in $\beta$. In each endpoint
block the edge complex is therefore the closed star of the corresponding
face, with the extra simplex already adjoined. These embeddings are full.
The construction is transported equivariantly to every coset vertex
and edge. Lemma~\ref{T-lem:tree} gives a flag complex $X$.

Every block chamber has one vertex of each type in $I_A\cup I_B$.
Flag transitivity in the two input geometries identifies its chambers
with the appropriate cosets of $H$. On an edge, the same identification
is by $C/H$, so these labels agree under every gluing. Thus the chambers
of $X$ are naturally $G/H$: equality of chamber labels is preserved
along the unique path of blocks supporting the chamber. In particular
$G$ acts transitively on them and their stabilizers are conjugates of $H$.

We give the same argument for an arbitrary standard face of type $J$,
including the empty face. Its stabilizer in the $A$ block is
$D_A=A_{J\cap I_A}$, its stabilizer in the $B$ block is
$D_B=B_{J\cap I_B}$, and its stabilizer on the edge is
$D_C=C\cap D_A=C\cap D_B$. The equality uses the fixed compatible
identification, not merely an abstract isomorphism of these groups.
The copies of this face that are identified across edges are exactly
the orbit of the standard copy under $\langle D_A,D_B\rangle$.
To justify this assertion locally, an edge $aC$ at the $A$ vertex
contains the standard $J$ face precisely when its $P$ face $aP$ is
incident with the standard $J\cap I_A$ face. Their union is then a
flag. Input flag transitivity supplies $d\in D_A$ carrying the
standard $P$ face to $aP$, so $aC=dC$. Thus the supporting incident
edges are exactly $D_A C/C$. At the other endpoint they are exactly
$D_B C/C$. A path transporting the face crosses only these edge
copies, so its consecutive changes of representatives are in $D_A$ or $D_B$.
Conversely each such change fixes the face. The reduced word theorem
identifies this generated subgroup with $D_A*_{D_C}D_B$.

Taking a one-vertex face identifies its stabilizer with the stated $G_i$.
Thus vertices of $X$ are precisely the corresponding cosets. Two vertices
in a simplex have a common chamber and hence a common coset representative;
conversely a common representative places them in a translate of the
standard chamber. Their incidence is therefore exactly coset incidence.
Taking the whole standard $J$ face gives the displayed intersection
formula, because its stabilizer is the intersection of its vertex
stabilizers. Finally every simplex of $X$ is contained in a block chamber,
so all flags extend to chambers. This proves the geometric assertion.
\end{proof}

\begin{corollary}\label{T-cor:residues}
Under the hypotheses of Theorem~\ref{T-thm:amalgam}, the amalgam geometry
is residually connected if and only if both input geometries are.
It is firm, thin or thick if and only if both inputs have the respective
property. Here the group actions need not be faithful.
\end{corollary}
\begin{proof}
The residue of a standard face of type $J$ has group $G_J$ and
parabolics $G_{J\cup\{i\}}$, for $i\notin J$. The intersection formula
therefore describes it by the same construction, using the corresponding
residues of the two input geometries. If $P\subseteq J$, its $A$ factor
is already in $C$ and the amalgam collapses to the $B$ factor. The
residue is exactly the residue of $\beta$ at $J\cap I_B$. If
$Q\subseteq J$, the symmetric statement holds. In particular, fixing
the standard $Q$ face recovers $\alpha$, and fixing the standard $P$
face recovers $\beta$; an empty such face simply means the whole geometry.
This proves necessity of residual connectivity.

For sufficiency, take a residue with at least two types left. The cases
just described reduce to residues of the inputs. Otherwise both
$P\setminus J$ and $Q\setminus J$ are nonempty. In the tree construction
of this residue, each $A$ block is a cone on its nonempty remaining
$Q$ simplex, and each $B$ block is a cone on its nonempty remaining
$P$ simplex. Each block is connected. The edge complexes contain the
nonempty simplex on $(P\cup Q)\setminus J$, so adjacent blocks meet.
Their union over a connected tree is connected. This includes all
residues of rank at least two and proves residual connectivity.

For a panel, $J=I\setminus\{i\}$ necessarily contains $P$ or $Q$.
The panel is consequently exactly an input panel by the first paragraph.
If $i\in L$, the two possible descriptions agree under the fixed
identification in $C$. Thus the panel sizes are the original indices
$[A_{I_A\setminus\{i\}}:H]$ or
$[B_{I_B\setminus\{i\}}:H]$, with every original index represented.
Comparison with $1$, $2$ or $3$ gives the three assertions.
\end{proof}


\subsection*{Converse for a compatible amalgam}
The converse uses a single-subgroup embedding, which is valid without the false unrestricted intersection formula for pairs of special subgroups.
Let $I_A,I_B$ be arbitrary finite type sets, put
$L=I_A\cap I_B$, $P=I_A\setminus L$, $Q=I_B\setminus L$,
and suppose a fixed isomorphism identifies $C=A_P$ with $B_Q$
and identifies $A_{P\cup J}$ with $B_{Q\cup J}$ for every
$J\subseteq L$. No flag assumption is imposed initially.
Set $G=A*_C B$ and
$G_i=\langle A_{\{i\}\cap I_A},B_{\{i\}\cap I_B}\rangle$.

\begin{proposition}\label{F-prop:converse}
If $(G,(G_i))$ is flag-transitive, both input incidence systems are
flag-transitive. Together with Theorem~\ref{T-thm:amalgam}, this is an
if and only if statement.
\end{proposition}
\begin{proof}
For $q\in Q$, we have $G_q=\langle A,B_q\rangle$ and
$C\leq B_q$. Fix left transversals for $C$ in $A$ and $B$.
A reduced amalgam word belongs to $\langle A,B_q\rangle$
precisely when all its $B$ syllables belong to $B_q$. This criterion
is independent of the choices of representatives, since multiplying
a syllable by an element of $C$ preserves membership in $B_q$.
It follows from uniqueness of the amalgam normal form, or by choosing
the $B$ transversal to extend one for $C$ in $B_q$.
Intersecting over $q\in Q$, every such syllable must belong to
$\bigcap_{q\in Q}B_q=C$, so there are no nontrivial $B$
syllables. Hence $G_Q=A$. If $Q$ is empty then $C=B$ and
$G=A$, so the same equality holds. Symmetrically $G_P=B$.

For $i\in I_A$, the two generating subgroups of $G_i$ have the
same intersection with $C$, by the fixed compatibility assumption.
The reduced word theorem embeds their amalgam in $G$ and gives
$G_i\cap A=A_i$. This is the single-subgroup embedding statement;
it does not use the false unrestricted intersection formula for two
special amalgam subgroups. In the flag-transitive geometry of $G$,
the residue of its standard $Q$ face is therefore the coset geometry
with group $G_Q=A$ and parabolics $G_Q\cap G_i=A_i$,
$i\in I_A$. Every residue of a flag-transitive geometry is
flag-transitive. This proves the assertion for $\alpha$, and the
standard $P$ residue gives $\beta$. Empty private type sets cause
no exception: the corresponding residue is the whole geometry.
\end{proof}


\subsection{Free products and their panel indices}
\begin{lemma}\label{F-lem:free-helly}
Let $\mathcal P$ and $\mathcal Q$ be families of subgroups of $A$ and
$B$, respectively, such that their left cosets are Helly. Then the
left cosets of the subgroups $P*Q\leq A*B$, with
$P\in\mathcal P$ and $Q\in\mathcal Q$, are Helly. Moreover,
\[
 \bigcap_{j=1}^m(P_j*Q_j)
   =\left(\bigcap_{j=1}^mP_j\right)*
     \left(\bigcap_{j=1}^mQ_j\right).
\]
\end{lemma}
\begin{proof}
The intersection assertion follows from the unique reduced word in a
free product: a word belongs to $P*Q$ exactly when each of its
$A$ syllables belongs to $P$ and each of its $B$ syllables belongs to
$Q$. This also proves $(P*Q)\cap A=P$ and $(P*Q)\cap B=Q$.

Consider the tree with vertices $gA,gB$ and with an edge labelled $g$
joining them, for every $g\in A*B$. For $H=P*Q$, the union of the
edges labelled by elements of $H$ is a connected subtree. Indeed a
reduced word in $P$ and $Q$ gives a path of such edges from the
identity edge to its own edge. Conversely the inclusions of the two
factors make this the usual tree of $P*Q$ embedded in the ambient
tree. Denote the translate associated to the coset $gH$ by $T(gH)$.
Its edge labels are exactly the elements of $gH$.

Take finitely many pairwise intersecting cosets $g_jH_j$. Their
subtrees meet pairwise in edges and hence have a common vertex $v$,
by the finite Helly property for subtrees of a tree. Suppose
$v=xA$; the other case is identical. Each set
$g_jH_j\cap xA$ is nonempty and, after multiplication by $x^{-1}$,
is a left coset of $P_j$ in $A$. For two of our subtrees, choose an
edge whose label belongs to both original cosets. The path from that
edge to $v$ lies in both subtrees. Its edge incident with $v$ has a
single label belonging to both cosets and to $xA$. Thus these local
$P_j$ cosets meet pairwise. The assumed Helly property in $A$ gives a
common element, which belongs to every original coset. The tree
Helly fact follows, for example, by rooting the tree and choosing the
deepest of the nearest-to-root points of the finitely many subtrees.
\end{proof}

\begin{theorem}\label{F-thm:free}
Let $\alpha=(A,(A_i)_{i\in I_A})$ and
$\beta=(B,(B_i)_{i\in I_B})$ be coset incidence systems with disjoint
type sets. In $G=A*B$ put $G_i=A_i*B$ for $i\in I_A$ and
$G_i=A*B_i$ for $i\in I_B$. Then
\[
 G_J=A_{J\cap I_A}*B_{J\cap I_B}.
\]
The resulting system is a flag-transitive geometry if and only if
both input systems are flag-transitive. Under these conditions it is
residually connected if and only if both inputs are.
\end{theorem}
\begin{proof}
The displayed formula and forward flag assertion are Lemmas~\ref{F-lem:cosets} and~\ref{F-lem:free-helly}, applied to the families of all input parabolics. For necessity,
lift a flag of $\alpha$ to the corresponding $G_i$ cosets using the
same representatives in $A$. Pairwise intersections remain nonempty.
A common element in $G$, followed by the retraction $G\to A$ that
kills $B$, gives a common element of the original cosets, since
$G_i$ maps onto $A_i$ for $i\in I_A$. This proves flag transitivity
of $\alpha$ without any Borel assumption. The retraction to $B$
gives the other input.

For a flag-transitive geometry, connectedness of a residue of rank at
least two is equivalent to generation of its group by the maximal
parabolics in that residue. To see this criterion, chambers incident
along an element connect their representatives through the appropriate
parabolic coset; the component containing the standard chamber consists
exactly of the chambers represented by the generated subgroup. Flag
transitivity identifies every residue with its parabolic coset system.

Apply this criterion to a face of type $J$. If both
$I_A\setminus J$ and $I_B\setminus J$ are nonempty, a remaining
$A$ type has residue parabolic containing the whole factor
$B_{J\cap I_B}$, and a remaining $B$ type has one containing
$A_{J\cap I_A}$. These two factors generate $G_J$. If only $A$
types remain, there are at least two of them. The residue parabolics
generate $G_J=A_{J\cap I_A}*B_{I_B}$ exactly when their projections
generate $A_{J\cap I_A}$, which is the required residue condition
in $\alpha$. The case with only $B$ types is symmetric. Taking
$J=J_A\cup I_B$, or $J=I_A\cup J_B$, proves necessity for every
input residue as well.
\end{proof}

\begin{proposition}\label{F-prop:panel}
Put $H_A=A_{I_A}$ and $H_B=B_{I_B}$. For a type $i\in I_A$,
write $U=A_{I_A\setminus\{i\}}$ and $U_0=H_A$.
The panel size in the free-product geometry is
\[
 [U*H_B:U_0*H_B]=
 \begin{cases}
 1,&U=U_0,\\
 [U:U_0],&H_B=1,\\
 \infty,&U_0<U\text{ and }H_B\ne1.
 \end{cases}
\]
There is the symmetric formula for a type in $I_B$. Consequently the
free-product geometry is firm if and only if both inputs are firm.
If both type sets are nonempty, it is thin if and only if both inputs
are thin and $H_A=H_B=1$.
\end{proposition}
\begin{proof}
The Borel is $H_A*H_B$ by Theorem~\ref{F-thm:free}, giving the index.
The first two cases follow directly. For the last, choose
$u\in U\setminus U_0$ and $1\ne v\in H_B$. The cosets
$(uv)^n(U_0*H_B)$, $n\geq0$, are distinct: equality would put a
nonzero power of $uv$ in $U_0*H_B$, contrary to its reduced word.
This proves infinitude, including when $U$ or $H_B$ is infinite.
An input panel has size one precisely when $U=U_0$, which proves
firmness. Under thinness all input panels must be nontrivial. If
both type sets occur, the last case then forces both Borels to be
trivial; the second case recovers exactly the input panel sizes.
The converse is immediate from the same formulas.
\end{proof}

\section{HNN extensions and indexed parabolics}\label{sec:hnn}


Let $A$ be a group and let $\mathcal P$ be a family of its subgroups,
closed under finite intersections. We say that $\mathcal P$ has the
\emph{coset Helly property} if every finite family of pairwise intersecting
left cosets of members of $\mathcal P$ has nonempty total intersection.
We include $A$ in $\mathcal P$. By Lemma~\ref{F-lem:cosets}, all standard parabolics of a finite-rank flag-transitive coset geometry form such a family.

Fix $C,D\in\mathcal P$ and an isomorphism $\phi:C\longrightarrow D$.
Write
\[
 G=\langle A,t\mid t^{-1}ct=\phi(c)\ (c\in C)\rangle.
\]
A subgroup $P\in\mathcal P$ is called compatible if
$\phi(P\cap C)=P\cap D$. Put $\widehat P=\langle P,t\rangle$.
Britton's lemma embeds the HNN extension of $P$ along these intersections
in $G$ and gives $\widehat P\cap A=P$. These are retained classical
normal-form facts. Unlike an arbitrary family of special subgroups, the
family below keeps its local coset Helly hypothesis.

\begin{theorem}\label{H-thm:helly}
If $\mathcal P$ has the coset Helly property, then the family consisting
of all $P\in\mathcal P$ and all $\widehat P$ for compatible
$P\in\mathcal P$ has the coset Helly property in $G$.
For compatible $P,Q$,
\[
 \widehat P\cap\widehat Q=\widehat{P\cap Q},\qquad
 R\cap\widehat P=R\cap P\quad(R\in\mathcal P).
\]
\end{theorem}
\begin{proof}
Use the Bass--Serre tree $T$ of this HNN extension. Its vertices are
$gA$ and its oriented edges are $gC$, with endpoints $gA$ and $gtA$.
The latter endpoint is independent of the representative because
$ct=t\phi(c)$. At a vertex $xA$, outgoing edge directions are
parametrized by left cosets of $C$ in $A$, and incoming directions by
left cosets of $D$ in $A$.

For compatible $P$, the embedded Bass--Serre tree of $\widehat P$ has
vertices $hA$ and edges $hC$, $h\in\widehat P$. Its injectivity follows
from the same reduced word theorem, since the associated intersections
are exactly $P\cap C$ and $P\cap D$. Denote this connected subtree by
$T_P$. For a base subgroup $R\leq A$ instead use the singleton subtree
$\{A\}$. Attach to a coset $g\widehat P$ the subtree $gT_P$ and to a
coset $gR$ the vertex $gA$. These assignments are independent of the
chosen coset representative.

If a vertex $v=xA$ belongs to such a support subtree, the corresponding
coset meets $xA$. After multiplication by $x^{-1}$ this intersection is
a left coset in $A$ of $P$, or of $R$ in the base-subgroup case.
For example, choose $y\in g\widehat P\cap xA$; the intersection is
$y(\widehat P\cap A)=yP$. An incident edge belongs to $gT_P$ precisely
when this local $P$ coset meets the local $C$ or $D$ coset describing
the direction of that edge. This follows directly from the vertex and
edge descriptions of the embedded tree.

We need one transport observation. Suppose two global cosets meet, and
a vertex $v$ lies in both support subtrees. Choose a common element $z$
of the cosets. Both support subtrees contain $zA$, so they contain the
whole finite path from $zA$ to $v$. Start with the common representative
$z$ at its first vertex. At the next edge, the two local subgroup cosets
intersect each other and each meets the local associated-subgroup coset
of that edge. Local coset Helly supplies a common element of all three.
For an outgoing edge multiply this representative on the right by $t$;
for an incoming edge multiply it by $t^{-1}$. The result lies at the
next vertex and still belongs to both global cosets, because both
subgroups contain $t$. If either subgroup was a base subgroup, its
support was a singleton and the path has length zero, so no multiplication
is needed. Repeating the step gives a common representative in $xA$.

Now take finitely many pairwise intersecting global cosets. Their support
subtrees intersect pairwise. The elementary Helly property for subtrees
gives a common vertex $v=xA$. By the transport observation every pair
of local fibers at $v$ intersects. These fibers are cosets of members
of $\mathcal P$, so local coset Helly supplies an element in all of
them. Multiplying by $x$ proves the first assertion.

For the intersection formula, first observe that $P\cap Q$ is compatible:
an isomorphism commutes with intersections. Thus
$K=\widehat{P\cap Q}$ is contained in both groups. Take
$g\in\widehat P\cap\widehat Q$ and follow the finite path from $A$
to $gA$. Both embedded subtrees contain it. We show inductively that
every vertex on this path is $kA$ for some $k\in K$.
This holds initially with $k=1$. Translate a current vertex $kA$ and
its next edge by $k^{-1}$. If that edge is outgoing, its local coset
$aC$ meets both $P$ and $Q$; if incoming, its local coset $aD$ does.
Since $P\cap Q$ contains $1$, the three local cosets meet pairwise.
The local Helly property supplies $p\in P\cap Q$ in the relevant
coset. The next vertex is therefore $kp tA$ or $kp t^{-1}A$,
respectively, and its representative is again in $K$.
At the end $gA=kA$, and hence
$k^{-1}g\in A\cap\widehat P\cap\widehat Q=P\cap Q\leq K$.
This proves the first equality. The second follows immediately from
$\widehat P\cap A=P$.
\end{proof}

\begin{corollary}\label{H-cor:geometry}
Let $(A,(A_i)_{i\in I})$ be a flag-transitive coset geometry of finite
rank. Let $C,D$ be standard parabolics, and let $\phi:C\to D$ be an
isomorphism. Choose a partition $\mathcal S$ of $I$ such that every
$A_s=\bigcap_{i\in s}A_i$, $s\in\mathcal S$, is compatible, and choose
$F\subseteq I$. In the HNN extension take maximal parabolics
\[
 G_s=\langle A_s,t\rangle\quad(s\in\mathcal S),\qquad
 G_*=A_F.
\]
This is a flag-transitive coset geometry. For a set $J$ of its types,
write $E=\bigcup_{s\in J\cap\mathcal S}s$. Then
\[
 G_J=
 \begin{cases}
   \langle A_E,t\rangle,&*\notin J,\\
   A_{E\cup F},&*\in J.
 \end{cases}
\]
Its chamber stabilizer is $A_I$.
\end{corollary}
\begin{proof}
Take $\mathcal P$ to be all standard parabolics of the input, including
$A$. These have the local coset Helly property explained above. Theorem
\ref{H-thm:helly} applies to all the displayed maximal parabolics. A flag
is a finite family of pairwise intersecting such cosets, so it has a
common representative. It is therefore a translate of a standard flag
and extends to a standard chamber. The intersection formulas follow
by induction from the theorem, with the usual empty-intersection convention.
For the complete type set $E=I$, so the Borel is $A_I$.
\end{proof}

\begin{corollary}\label{H-cor:rc}
Suppose the input in Corollary~\ref{H-cor:geometry} is residually connected.
Let $C=A_J$, $D=A_K$, and suppose a specified bijection
$f:I\setminus J\longrightarrow I\setminus K$ satisfies
\[
 \phi(A_{J\cup X})=A_{K\cup f(X)}\qquad(X\subseteq I\setminus J).
\]
Take $\mathcal S$ to be the connected classes of the arrows $i\to f(i)$
on $I$, and put $F=J\setminus K$. Then the resulting HNN geometry is
residually connected. Its $*$ panels are infinite. Its panel of type
$s\in\mathcal S$ has size
\[
 [A_{I\setminus(s\setminus F)}:A_I].
\]
In particular it is firm when the input is firm.
\end{corollary}
\begin{proof}
Write $H=A_I$ and $R_i=A_{I\setminus\{i\}}$.
The displayed compatibility gives $\phi(H)=H$ and
$\phi(R_i)=R_{f(i)}$ for $i\notin J$.
Each class $s$ satisfies $f(s\setminus J)=s\setminus K$.
Consequently $\phi(A_s\cap C)=A_s\cap D$, so the previous corollary
applies. Its minimal parabolics are
\[
 M_s=A_{I\setminus(s\setminus F)},\qquad M_*=\langle H,t\rangle.
\]

Residual connectivity of the input gives the elementary generation rule
\[
 A_{I\setminus E}=\langle H,R_i\ (i\in E)\rangle
 \quad(E\subseteq I).
\]
This is Lemma~\ref{O-lem:parabolic-join}, including the empty support; the Borel is not dropped.

Let $W$ be a subset of the new type set $\mathcal S\cup\{*\}$, and
put $E=\bigcup_{s\in W\cap\mathcal S}s$. If $*\notin W$, then
\[
 \langle H,M_s\ (s\in W)\rangle
 =A_{I\setminus(E\setminus F)}=G_{(\mathcal S\cup\{*\})\setminus W}.
\]
If $*\in W$, this generated group also contains $t$.
Every $i\in E\cap F$ is in the range but not the domain of $f$.
It therefore has a predecessor $j\in E\setminus F$ with $f(j)=i$.
Conjugation by $t$ brings $R_i$ into the group from $R_j$.
Thus in this case the generated group is
$\langle H,t,R_i\ (i\in E)\rangle=\langle A_{I\setminus E},t\rangle$,
again exactly $G_{(\mathcal S\cup\{*\})\setminus W}$.

We have proved the same generation rule for every new standard parabolic.
If at least two types remain, each of its minimal generators belongs to
one of the parabolics obtained by deleting a different remaining type.
Those parabolics therefore generate it. The residual generation criterion
proves residual connectivity.

Since $\phi(H)=H$, the subgroup $\langle H,t\rangle$ is an extension
of $H$ by the infinite cyclic group generated by the stable letter.
Its index over $H$ is infinite. The other panel indices are the stated
indices of $M_s$ over $H$. Every class has a point outside $F$: a point
in $F$ has a predecessor outside $F$. If the input is firm, the group
$M_s$ therefore contains some $R_i$ strictly larger than $H$.
All these panels have at least two elements, as required.
\end{proof}


\subsection{Firmness recovers the missing indices}


Let $\alpha=(A,(A_i)_{i\in I})$ be a coset geometry of finite rank.
Write $A_J=\bigcap_{j\in J}A_j$, $A_\varnothing=A$,
$H=A_I$, and $R_i=A_{I\setminus\{i\}}$.
The following criterion turns an isomorphism of unlabelled parabolic lattices into the indexed compatibility needed above. Firmness is an actual hypothesis.

\begin{lemma}\label{I-lem:boolean}
Suppose $R_i>H$ for every $i\in I$. Then
\[
 A_J\leq A_K\quad\Longleftrightarrow\quad K\subseteq J.
\]
Consequently the lattice of standard parabolics, ordered by inclusion,
is the reverse Boolean lattice on $I$.
\end{lemma}
\begin{proof}
If $K\subseteq J$, the inclusion follows from the definition of an
intersection. Conversely, suppose $k\in K\setminus J$. Then
$R_k\leq A_J$, because every member of $J$ differs from $k$.
If $A_J\leq A_K$, it follows that $R_k\leq A_k$, whence
$R_k=R_k\cap A_k=H$, a contradiction. This proves the equivalence,
including injectivity of $J\mapsto A_J$.
\end{proof}

\begin{corollary}\label{I-cor:firm}
Suppose $\alpha$ is flag-transitive and firm. Let
$\phi:A_J\longrightarrow A_K$ be an isomorphism inducing an isomorphism
of the unlabelled lattices of standard parabolics contained in its domain
and codomain. Then there is a unique bijection
$f:I\setminus J\longrightarrow I\setminus K$ such that
\[
 \phi(A_{J\cup X})=A_{K\cup f(X)}
 \qquad(X\subseteq I\setminus J).
\]
Thus the indexed hypothesis of Corollary~\ref{H-cor:rc} follows from the source's
admissibility condition for firm flag-transitive inputs.
\end{corollary}
\begin{proof}
Flag transitivity identifies the chambers on a standard panel of missing
type $i$ with $R_i/H$. Firmness therefore implies $[R_i:H]\geq2$.
Lemma~\ref{I-lem:boolean} applies. The standard parabolics contained in
$A_J$ are precisely the $A_{J\cup X}$, and their coatoms are
$A_{J\cup\{i\}}$ for $i\notin J$. The isomorphism permutes coatoms
bijectively, giving $f$. Since it preserves intersections, intersecting
these coatom equalities proves the formula for every $X$, including
$X=\varnothing$.

For completeness, form the equivalence classes $s$ of the arrows
$i\mapsto f(i)$ on $I$. Bijection of the domain and range gives
$f(s\setminus J)=s\setminus K$. Hence
\[
 \phi(A_s\cap A_J)
 =\bigcap_{i\in s\setminus J}\phi(A_{J\cup\{i\}})
 =\bigcap_{j\in s\setminus K}A_{K\cup\{j\}}
 =A_s\cap A_K.
\]
Both displayed intersections of images are taken inside $A_K$; if
their index sets are empty they mean $A_K$, in agreement with
$\phi(A_J)=A_K$. Every $A_s$ is therefore compatible. These are precisely the
subgroups required by the local coset Helly theorem. Its flag-transitivity
conclusion follows; its residual-connectivity conclusion also requires
residual connectivity of the input, as stated there.
\end{proof}


\subsection{Residually connected inputs with universal types}
Let $\alpha=(A,(A_i)_{i\in I})$ be a finite-rank flag-transitive,
residually connected coset geometry. Write $A_J=\bigcap_{i\in J}A_i$,
$H=A_I$, and $R_i=A_{I\setminus\{i\}}$.
The Borel-preserving residual generation formula is
\begin{equation}\label{R-eq:gen}
 A_J=\langle H,R_i:i\notin J\rangle.
\end{equation}
This is the same Borel-preserving formula proved in Lemma~\ref{O-lem:parabolic-join}.

\begin{lemma}\label{R-lem:active}
Put $I_0=\{i:A_i=A\}$ and $I_+=I\setminus I_0$.
Then $A_i=A$ if and only if $R_i=H$. Deleting the universal types
$I_0$ leaves a firm flag-transitive, residually connected geometry
on $I_+$ with the same unlabelled family of standard parabolics and
the same Borel $H$. The original incidence geometry is its join
with one singleton for each type in $I_0$.
\end{lemma}
\begin{proof}
If $A_i=A$, then $H=A_i\cap R_i=R_i$. Conversely, if $R_i=H$,
formula (\ref{R-eq:gen}) for $J=\varnothing$ and $J=\{i\}$ gives
$A=A_i$. A type with maximal parabolic $A$ has exactly one element,
incident with every other element, proving the join description.
Dropping these factors leaves every standard intersection unchanged.
Flags still have common representatives. Every residue on at least
two active types is a residue of the original geometry after fixing
the deleted universal elements, so it remains connected. Each active
panel has index $[R_i:H]>1$, proving firmness. The empty active
geometry is interpreted as a single empty chamber.
\end{proof}

For later comparison we also make the firm branch explicit without
assuming residual connectedness. For a firm flag-transitive input
and a source-admissible $\phi:A_J\to A_K$,
Lemma~\ref{I-lem:boolean} gives the indexed partial
bijection required by Theorem~\ref{H-thm:helly}. For its classes
$s$ and $F=J\setminus K$, that theorem's intersection formula gives
the Borel $H$, the stable minimal parabolic $\langle H,t\rangle$,
and the other minimal parabolics
\[
 M_s=A_{I\setminus(s\setminus F)}.
\]
The first has infinite index over $H$, since the stable letter has
infinite image in $\ZZ$ and normalizes $H$. Every class has an
element outside $F$: a point in $F$ is in the range but not the
domain of the partial bijection, and its predecessor is outside
$J$, hence outside $F$. Consequently $M_s$ contains some
$R_i>H$. All panels are firm. This proves the full firm-input branch
of source Theorem 3.35, including part (c), independently of the
residual-connectivity assertion below.

Suppose $C=A_J,D=A_K$, and $\phi:C\to D$ satisfies the source's
Definition 3.30: it induces an isomorphism of the unlabelled families
of standard parabolics lying in these subgroups. The source forms
classes of the arrows $i\to j$, with $i\notin J,j\notin K$, for
which $\phi(A_{J\cup\{i\}})=A_{K\cup\{j\}}$. It uses
$F=J\setminus K$, the maximal subgroups $\langle A_s,t\rangle$
for these connected classes $s$, and the stable-type subgroup $A_F$
in the HNN extension $\langle A,t\mid t^{-1}ct=\phi(c)\rangle$.

\begin{theorem}\label{R-thm:rc}
With these literal choices, every flag-transitive, residually connected
input gives a flag-transitive, residually connected output. Thus the
residually connected branch of Theorem 3.35 remains valid without
firmness. Active output panels have the indices in the indexed HNN theorem; each additional universal class has one element
in its panel.
\end{theorem}
\begin{proof}
Apply Lemma~\ref{R-lem:active}. On the active types the firm geometry
has the reverse Boolean parabolic lattice proved in Lemma~\ref{I-lem:boolean}. Thus the isomorphism $\phi$ gives a
bijection
\[
 f:I_+\setminus J\longrightarrow I_+\setminus K
\]
with
$\phi(A_{J\cup X})=A_{K\cup f(X)}$ for every
$X\subseteq I_+\setminus J$. Here intersections with $J,K$
can ignore their inactive entries. This follows because the coatoms
of the Boolean interval inside $A_J$ are exactly the subgroups
$A_{J\cup\{i\}}$ for its missing active types; an isomorphism
permutes those coatoms and commutes with all their intersections.
The minimal subgroup $H$ is also preserved.

For an active $i\notin J$, the subgroup $A_{J\cup\{i\}}$ is
proper in $A_J$, whereas for an inactive $i$ it equals $A_J$.
The same distinction holds on the $K$ side. Hence a source arrow
never connects an active type to an inactive one. On active types
the connected classes are exactly those of the partial bijection
$f$, with isolated active types included. The stable parabolic is
$A_F=A_{F\cap I_+}$. Thus all active-class parabolics and the
stable parabolic coincide exactly with the indexed HNN
construction of Corollary~\ref{H-cor:rc} on the firm active geometry. Its flag transitivity and
residual connectedness follow from that theorem.

An inactive class $s$ consists only of types with $A_i=A$, so
$A_s=A$ and $\langle A_s,t\rangle$ is the entire HNN group.
Its coset type is a singleton universal element. Adding any number
of these types preserves flag transitivity and residual connectedness:
a residue with at least two active types uses the active conclusion,
a residue with one active type and at least one universal type has
a connected star, and a residue consisting of at least two universal
types is a complete graph. The same observation covers the case
$I_+=\varnothing$. In that case $C=D=A$ and the stable-type cosets
of $A$ in $A\rtimes\ZZ$ are the infinite rank-one part; all other
types are universal. This proves the assertion and panel description.
\end{proof}

\section{Semidirect products and orbit intersections}\label{sec:semidirect}

Let $\alpha=(A,(A_i)_{i\in S})$ and
$\beta=(B,(B_j)_{j\in T})$ be flag-transitive coset geometries of
finite rank. Write $H_A=A_S$ and $H_B=B_T$. Suppose that $B$ acts on
$A$ by automorphisms and that an indexed action on $S$ is specified,
with $b(A_i)=A_{bi}$. This specification also covers coincident
parabolics: an action on an unlabelled set of subgroups alone does not
automatically specify an action on their repeated indices. Let
$\mathcal O$ be the orbit partition of $S$.

For an orbit $s\in\mathcal O$, put $A_s=\bigcap_{i\in s}A_i$.
In $G=A\rtimes B$ take the indexed parabolics
\[
 P_s=A_s\rtimes B\quad(s\in\mathcal O),\qquad
 P_j=A\rtimes B_j\quad(j\in T).
\]
The orbit types here use the whole orbit intersection. The construction below instead uses orbits of parabolic subgroups; these are different incidence structures with different panel formulas.

\begin{proposition}\label{J-prop:join}
Let $\overline\alpha=(A,(A_s)_{s\in\mathcal O})$. The coset
incidence system $(G,(P_i))$ is the join
$\overline\alpha*\beta$: the two component incidences are unchanged
and every element of one component is incident with every element of
the other. Both $\overline\alpha$ and the join are flag-transitive
geometries. For $J=J_O\sqcup J_T$, put
$E=\bigcup_{s\in J_O}s$. Then
\[
 G_J=A_E\rtimes B_{J_T},\qquad G_{\mathcal O\sqcup T}=H_A\rtimes H_B.
\]
\end{proposition}
\begin{proof}
Every $A_s$ is $B$-invariant. In semidirect coordinates,
\[
 (a,b)P_s=(aA_s)\times B,\qquad
 (a,b)P_j=A\times(bB_j).
\]
Here the second equality uses that $b$ acts bijectively on $A$.
Thus an orbit-type element is determined by the coset $aA_s$, a
$T$-type element by $bB_j$, and their incidences are exactly those
stated. The parabolic formula follows by intersecting the two
coordinates; this is ordinary intersection of Cartesian products.

The cosets of the $A_s$ belong to the family of cosets of standard
parabolics of $\alpha$. That family is Helly by Lemma~\ref{F-lem:cosets}. Hence $A$ is flag-transitive on $\overline\alpha$, and every
flag there extends to a chamber. The same is assumed for $\beta$.
For a flag of the join, choose $b\in B$ carrying its $\beta$ part
to the standard flag. This acts on its $\overline\alpha$ part;
then an element of $A$ carries that part to the standard flag while
fixing the $\beta$ part. Thus $G$ is flag-transitive. Alternatively,
a common $B$ coordinate and a common $A$ coordinate in the displayed
cosets give a common group representative directly.
\end{proof}

\begin{proposition}\label{J-prop:residues}
If both inputs are residually connected, so is the semidirect join.
Its panel indices are
\[
 [A_{S\setminus s}:H_A]\quad(s\in\mathcal O),\qquad
 [B_{T\setminus\{j\}}:H_B]\quad(j\in T).
\]
It is firm when both inputs are firm. If both inputs are thin, the
join is thin exactly when every orbit in $S$ is a singleton.
\end{proposition}
\begin{proof}
For clarity we first prove residual connectivity of
$\overline\alpha$. Put $R_i=A_{S\setminus\{i\}}$.
Residual connectivity of the finite-rank flag-transitive input gives
\[
 A_{S\setminus E}=\langle H_A,R_i:i\in E\rangle
 \qquad(E\subseteq S).
\]
This is Lemma~\ref{O-lem:parabolic-join}. Take a coarse
residue with at least two orbit types remaining, whose union is $E$.
For every $i\in E$ choose a remaining orbit $s$ not containing $i$.
Then $R_i\leq A_{(S\setminus E)\cup s}$, one of the maximal
parabolics of that residue. Those parabolics also contain $H_A$ and
therefore generate its whole group $A_{S\setminus E}$. This proves
residual connectivity of the coarse geometry.

A residue of a join is the join of the corresponding component
residues. If both components have a type remaining, the incidence
graph is connected because all vertices across the two components
are incident. If just one component has types remaining, residual
connectivity is that of the corresponding input residue. This proves
the first assertion, including empty component type sets.

The panel formulas follow by omitting just one type from the
parabolic formula in Proposition~\ref{J-prop:join}. For an orbit $s$,
choose $i\in s$. Firmness of $\alpha$ gives $H_A<R_i$, and
$R_i\leq A_{S\setminus s}$, so its panel index is at least two.
The other panels have exactly their $\beta$ sizes.

Suppose both inputs are thin. Singleton orbit indices are two. If
$s$ contains distinct $i,k$, choose $x\in R_i\setminus H_A$ and
$y\in R_k\setminus H_A$. We have $R_i\cap R_k=H_A$ by their
definitions. The four left cosets
\[
 H_A,\quad xH_A,\quad yH_A,\quad xyH_A
\]
are distinct. For example, $xyH_A=xH_A$ would put $y$ in $H_A$;
$xyH_A=yH_A$ would put $x\in yH_Ay^{-1}\leq R_k$, a
contradiction; and $xyH_A=H_A$ would put $x\in R_k$.
Thus $[A_{S\setminus s}:H_A]\geq4$, ruling out thinness.
The argument uses only subgroup intersections and does not assume
that $H_A$ is normal in all of $A$.
\end{proof}


\subsection{A general orbit construction}
Let $\alpha=(A,(A_x)_{x\in S})$ and
$\beta=(B,(B_i)_{i\in I})$ be flag-transitive coset geometries.
Suppose a specified homomorphism $B\to\Aut(A)$ and an action of $B$
on the indexed set $S$ satisfy $b(A_x)=A_{bx}$.
The indexed action is part of the data, even if different maximal
parabolics happen to coincide. Let $\mathcal K$ be the set of orbits
on $S$, choose $x_L\in L$ for each $L\in\mathcal K$, and assume
\begin{equation}\label{O-eq:general-orbits}
 B_Hx_L\cap B_Kx_L=B_{H\cup K}x_L
 \quad(H,K\subseteq I,\ L\in\mathcal K).
\end{equation}
For $T=I\sqcup\mathcal K$, define subgroups of $G=A\rtimes B$ by
\begin{equation}\label{O-eq:general-maximal}
 G_i=A_{S\setminus\bigcup_L B_ix_L}\rtimes B_i\quad(i\in I),
 \qquad G_L=A_L\rtimes B\quad(L\in\mathcal K).
\end{equation}
Every indicated subset of $S$ is invariant under its second factor,
so these are indeed subgroups. For $J\subseteq T$ write
$J_B=J\cap I$, $J_A=J\cap\mathcal K$, and set
\[
 E(J)=S\setminus\bigcup_{L\in\mathcal K\setminus J_A}B_{J_B}x_L.
\]

\begin{theorem}[Orbit construction with a Borel subgroup]\label{O-thm:general-twist}
The coset system in (\ref{O-eq:general-maximal}) is a flag-transitive
geometry. All its standard parabolics, including its Borel subgroup, are
\begin{equation}\label{O-eq:general-parabolics}
 G_J=A_{E(J)}\rtimes B_{J_B},\qquad G_T=A_S\rtimes B_I.
\end{equation}
It is finite if and only if both input geometries are finite. It is
residually connected if both inputs are residually connected, and firm
if both inputs are firm. Its numbers of chambers through panels are
\begin{align}
 [G_{T\setminus\{i\}}:G_T]&=[B_{I\setminus\{i\}}:B_I]
       &&(i\in I),\label{O-eq:old-panel}\\
 [G_{T\setminus\{L\}}:G_T]&=[A_{S\setminus B_Ix_L}:A_S]
       &&(L\in\mathcal K).\label{O-eq:new-panel}
\end{align}
In particular, if both inputs are thin and $B_I$ fixes each chosen
$x_L$, the new geometry is thin. Faithfulness of either input is not
needed for this sufficient condition.
\end{theorem}
\begin{proof}
Intersections of subgroups $A_E\rtimes B_H$ are computed coordinate by
coordinate. The intersection of their $A$-parts is obtained by taking
the union of their subscripts. Different $B$-orbits are disjoint, and
iteration of (\ref{O-eq:general-orbits}) says that the intersection of
the old supports is the support for $B_{J_B}$. Intersecting with the
orbit parabolics deletes precisely the orbits in $J_A$. Taking
complements proves (\ref{O-eq:general-parabolics}). It uses intersections
of maximal parabolics, not an unjustified formula for generated supports.

To prove flag transitivity, project a flag of $G$-cosets to $B$.
Old types give pairwise-intersecting maximal-parabolic cosets; orbit
types project to all of $B$. Coset intersection from the first section
gives a common $b\in B$. After left translation by $(1,b^{-1})$,
each coset has a representative $(a,1)$. Such a coset is the Cartesian
product of a coset of a standard parabolic of $A$ and the corresponding
subgroup of $B$. Pairwise incidence therefore gives a pairwise-intersecting
family of standard-parabolic $A$-cosets. Their common representative,
again by the coset intersection lemma, realizes the entire flag.
Thus every flag is an image of a standard flag and extends to a chamber.
This step by itself does not require the orbit identity; that identity
is needed for the exact parabolic description and the later conclusions.

Suppose both inputs are residually connected, and fix $J$ with at
least two missing types. We prove generation of $G_J$ by the next
parabolics. If no orbit type is missing, all their $A$-parts are $A_S$;
(\ref{O-eq:rc-borel}) for $B$ gives the required generation. If no old
type is missing, choose two distinct missing orbits $L,M$. The two
corresponding next parabolics both contain $B_I$. Their $A$-subscripts
are $E(J)\cup L$ and $E(J)\cup M$, whose intersection is $E(J)$.
Lemma~\ref{O-lem:parabolic-join} gives the entire $A$-part.

In the remaining case choose a missing old type $i$ and a missing
orbit type $L$. The parabolic $G_{J\cup\{L\}}$ contains the full
$B_{J_B}$. Write $E'=E(J\cup\{i\})$. The group generated by these
two next parabolics therefore contains $A_{bE'}$ for every
$b\in B_{J_B}$. In each remaining orbit, $B_{J_B\cup\{i\}}x_M$
is nonempty and contains $x_M$. Its $B_{J_B}$-saturation is exactly
$B_{J_B}x_M$. Consequently
\[
 \bigcap_{b\in B_{J_B}}bE'=E(J).
\]
There are only finitely many distinct subsets of the finite set $S$,
so Lemma~\ref{O-lem:parabolic-join} applies directly to these conjugate
parabolics. They generate $A_{E(J)}$, proving the remaining case.
The residue identification and connectivity criterion already proved
now establish residual connectedness in every residue.

Putting $J=T\setminus\{t\}$ in
(\ref{O-eq:general-parabolics}) gives the two panel formulas.
For an old type the $A$-parts agree; for an orbit type the $B$-parts
agree. These observations prove the stated indices even when the
Borel subgroups are nontrivial. If the inputs are firm, the old indices
are at least two. The orbit $B_Ix_L$ contains $x_L$, so
$A_{S\setminus\{x_L\}}\leq A_{S\setminus B_Ix_L}$; its index over
$A_S$ is at least two by firmness of $\alpha$. This proves firmness.
If the inputs are thin and $B_Ix_L=\{x_L\}$, all the displayed
indices are two, proving the sufficient thinness assertion.

Finally the number of chambers is the index of $G_T$. The elementary
coset index formula gives
\[
 [G:G_T]=[A:A_S][B:B_I],
\]
with the equality interpreted as a product of cardinalities.
For a finite number of types, a geometry has finitely many elements
if and only if it has finitely many chambers: every element extends
to a chamber and a chamber is a tuple of one element of each type.
The same argument applies to each input. Thus the product is finite
exactly when both input geometries are finite. No finiteness of the
possibly unfaithful groups $A$ and $B$ was assumed.
\end{proof}

\begin{corollary}[Faithful regular inputs]\label{O-thm:twist}
For faithful regular hypertopes $\alpha,\beta$, the construction is a
regular hypertope with full type-preserving automorphism group $A\rtimes B$.
Writing $A(R)=\langle a_x:x\in R\rangle$, its parabolics are
\begin{equation}\label{O-eq:regular-support}
 G_J=A\!\left(\bigcup_{L\notin J_A}B_{J_B}x_L\right)\rtimes B_{J_B}.
\end{equation}
Its distinguished involutions are the $b_i$ and the $a_{x_L}$.
\end{corollary}
\begin{proof}
The faithful regular model has $A_S=B_I=1$ and
$A_E=A(S\setminus E)$, as proved in the first section. The preceding
theorem supplies flag transitivity, residual connectedness and thinness.
Its Borel subgroup is trivial, and the panel groups are precisely the
cyclic groups on the indicated involutions. Chamber propagation shows
that the displayed regular action is the full type-preserving group.
\end{proof}
This recovers the general finite-type twisting theorem and its regular
specialization in Piedade--Tranchida~\cite[Theorem 3.16 and Corollary 3.17]{N006}.
The exact panel formula also explains why passing to a faithful quotient
of $B$ cannot silently preserve an arbitrary action on the types of
$\alpha$: its Borel subgroup must act trivially for that action to descend.

\section{Invariant quotients and orbital presentations}
\begin{lemma}[Quotients]\label{O-lem:quotient}
Use the faithful regular input geometries and specified indexed action of Corollary~\ref{O-thm:twist}. If $M\leq A$ is invariant under $B$, its normal closure $N$ in $A$ is invariant. The action descends to $A/N$ and permutes the images of the standard maximal parabolics with the same indexed type action. If their coset geometry is a regular hypertope, the corresponding twisting is again a regular hypertope.
\end{lemma}
\begin{proof}
The normal closure is generated by $a m a^{-1}$ with $a\in A,m\in M$. An automorphism from $B$ carries each such generator to another of the same form, and its inverse gives equality. The quotient action and parabolic permutation follow. Orbit identity (\ref{O-eq:general-orbits}) concerns only the indexed types and is unchanged.

For completeness, the regular-geometry assumption can be interpreted without silently assuming a faithful quotient action. Put $\bar A=A/N$, $C_x=A_xN/N$, and $C_E=\bigcap_{x\in E}C_x$, with $C_\varnothing=\bar A$. The source's quotient twisting uses intersections of these maximal parabolics:
\begin{equation}\label{O-eq:quotientparabolics}
\bar G_i=C_{S\setminus\bigcup_L B_ix_L}\rtimes B_i\quad(i\in I),
\qquad \bar G_L=C_L\rtimes B\quad(L\in\mathcal K).
\end{equation}
These are not defined by substituting the images of generated-support subgroups into (\ref{O-eq:regular-support}).

Each image of $a_x$ fixes the base flag with type $x$ deleted. It cannot belong to $C_x$: otherwise $C_x$, already containing the other generators, would be all of $\bar A$, giving only one element of type $x$, contrary to thinness. It therefore performs the unique adjacency of that type. Chamber-graph connectivity makes $\bar A$ flag-transitive. Its chamber stabilizer $H=C_S$ fixes every chamber by thinness, so is its normal action kernel. It is $B$-invariant because it is the intersection of the permuted $C_x$.

Every $C_E$ contains $H$, and intersections commute with passage to $\bar A/H$. In that faithful regular model formula (\ref{O-eq:parabolic}) therefore gives
\[
C_E/H=\bigcap_{x\in E}(C_x/H)
=\langle a_xH:x\in S\setminus E\rangle.
\]
Consequently every subgroup in (\ref{O-eq:quotientparabolics}) is the full inverse image of the corresponding subgroup from Corollary~\ref{O-thm:twist} for $\bar A/H$. Apply that theorem. Passage through the normal subgroup $H\rtimes1$ preserves all cosets and their intersections, so the original quotient twisting has the same regular geometry. For a transitive type orbit $L=S$, its new parabolic is $H\rtimes B$, not necessarily $B$. The quotient group $\bar A$ itself is not asserted to be the faithful automorphism group.
\end{proof}

\begin{proposition}[Orbital presentation]\label{O-prop:orbital}
Let $B$ act transitively on a finite set $X$ with base point $x_0$ and stabilizer $P$. Give each orbit of unordered distinct pairs a Coxeter label. Let $Y$ be a finite set fixed by $B$, with fixed labels within $Y$ and one label $m_y$ between $y$ and every element of $X$. Put $A=W(Y\sqcup X)$. Then $A\rtimes B$ has the presentation consisting of a presentation for $B$, one involution $a=a_{x_0}$, the involutions $c_y$ for $y\in Y$, and:
\begin{enumerate}
\item $[a,p]=1$ for generators $p$ of $P$, and $[c_y,b]=1$ for generators $b$ of $B$;
\item the prescribed relations among the $c_y$, and $(c_ya)^{m_y}=1$ when $m_y<\infty$;
\item $(a\tau a\tau^{-1})^{m}=1$ for one representative $\{x_0,\tau x_0\}$ of each unordered pair orbit of finite label $m$.
\end{enumerate}
The unordered orbits are indexed by the inverse-paired double cosets $P\tau P\cup P\tau^{-1}P$ with $\tau\notin P$.
\end{proposition}
\begin{proof}
In the proposed group define $a_{bx_0}=bab^{-1}$. The stabilizer relations make this independent of $b$. Double-coset multiplication changes the representative without changing its pair orbit, and inversion interchanges its endpoints. Conversely, a transporter of two unordered pairs either fixes the selected first endpoint or interchanges the endpoints, giving exactly the two double cosets displayed. Hence conjugates of the chosen relations give every Coxeter relation on $X$. Conjugating the relations with $c_y$, which is fixed by $B$, gives all the mixed Coxeter relations. These elements therefore define a homomorphism from $W(Y\sqcup X)\rtimes B$ onto the proposed group. The reverse map takes the proposed generators to their named elements in the semidirect product; all the relators hold there. Both composites fix the generating set, proving the isomorphism. This is a presentation of the full group, not merely a finite image.
\end{proof}

\begin{lemma}[Exact doubled label]\label{O-lem:double}
In $A\rtimes B$ let $a=a_x$ and let $b\in B$ be an involution with $bx\ne x$. If $a_xa_{bx}$ has order $m$ (possibly infinite), then $ab$ has order $2m$, with $2\infty=\infty$.
\end{lemma}
\begin{proof}
We have $(ab)^2=a_xa_{bx}$. Projection to $B$ sends $ab$ to the nontrivial involution $b$, so any power of $ab$ equal to $1$ has even exponent. For an even exponent $2k$ it is $1$ precisely when $(a_xa_{bx})^k=1$. This proves exactness in both the finite and infinite cases.
\end{proof}


\subsection{Indexed quotient actions without a geometric assumption}
For a group $A$ with a specified action of $B$, a $B$-invariant
subgroup $H$ has invariant normal closure $N$: every image of
$a^{-1}ha$ is a conjugate of an element of $H$, and applying the
inverse automorphism gives equality. For arbitrary $H$, the subgroup
$\langle b(h):b\in B,h\in H\rangle$ is invariant, by reindexing
$b$ after any further action. The statement applies independently of any geometric conclusion about the quotient.

The induced action on $A/N$ sends each image $A_iN/N$ to the image
indexed by $bi$. Keep this indexed action as part of the data even
if some subgroup images coincide. Every orbit-intersection condition
involving the acting group and named types is consequently unchanged.
The image subgroups are
$\widetilde A_i=A_iN/N$, rather than $A_i/N$, which need not be defined. No flag transitivity, intersection
property, firmness or residual connectedness is inferred for this
quotient. Those are separate hypotheses when applying twisting.

\section{Cyclic markings and exact presentations}\label{sec:cyclic}
\subsection{The precise input}
For a group generated by a finite named set $S$, write
$A(R)=\langle r:r\in R\rangle$ for $R\subseteq S$, including
$A(\varnothing)=1$. We call this a parabolic marking if
\begin{equation}\label{C-eq:ip}
 A(R)\cap A(Q)=A(R\cap Q)
\end{equation}
and the cosets of all $A(R)$ have the finite Helly property: every finite
pairwise-intersecting collection has a common element. The names denote
distinct nonidentity elements. The standard coset system has maximal
parabolics $A(S\setminus\{s\})$.

Here are the classical inputs retained for an arbitrary Artin--Tits
marking. Standard parabolic presentations embed, their intersections obey
(\ref{C-eq:ip}), and their cosets have this Helly property. These are the
standard-parabolic results invoked in N005 Theorems~2.11--2.12:
van der Lek's parabolic embedding and the intersection/product results
cited there from Godelle--Paris. In particular, each singleton Artin
parabolic is infinite cyclic. For a Coxeter marking the corresponding
facts are the classical Coxeter parabolic theorem, and singleton
parabolics have order two. We do not reprove those external
results or identify their proof costs with zero.
Precisely, the source locators are Bourbaki, \emph{Groupes et alg\`ebres
de Lie}, Chapters 4--6 (2006 edition), Chapter 4, \S1, Theorem 2
and Exercise \S1.1(b); H.~van der Lek, \emph{The homotopy type of
complex hyperplane complements} (1983), Chapter II, Theorem 4.13;
and E.~Godelle--L.~Paris, \emph{$K(\pi,1)$ and word problems for
infinite type Artin--Tits groups, and applications to virtual braid
groups}, Math.~Z.~272 (2012), Proposition 4.5. We retain these
interfaces as the primary source does, rather than claiming fresh
proof verification of the entire cited works.

Lemma~\ref{F-lem:cosets} identifies this Helly hypothesis with flag transitivity, once formula (\ref{C-eq:ip}) identifies generated supports with standard intersections.

\begin{lemma}\label{C-lem:product}
In a parabolic marking, for $H=A(R)$, $K=A(Q)$ and $L=A(P)$,
\[
 HK\cap HL=H(K\cap L).
\]
The complex of cosets of maximal parabolics, with simplices the families
having common intersection, is a flag complex. Its coset geometry is
residually connected; its panel of type $s$ has $|\langle s\rangle|$
chambers.
\end{lemma}
\begin{proof}
If $z\in HK\cap HL$, the left cosets $H,zK,zL$ meet pairwise:
the first meets each of the others and the latter two contain $z$.
Helly gives $h=zk=zl\in H$, with $k\in K,l\in L$; hence
$k=l\in K\cap L$ and $z\in H(K\cap L)$. The reverse inclusion
is immediate. The complex assertion is exactly Helly. Every standard
residue has acting group $A(R)$ and maximal parabolics
$A(R\setminus\{s\})$, $s\in R$. If $|R|\geq2$ they generate
$A(R)$: each named generator belongs to a child missing a different
name. Its incidence graph is therefore connected (successive
multiplications by parabolic elements give incidence paths). Flag
transitivity identifies every residue with a translate of one of these.
The Borel is $1$ by (\ref{C-eq:ip}), and the rank-one residue is
$\langle s\rangle/1$, proving the panel assertion.
\end{proof}

\subsection{One cyclic marking theorem}
Let $S=Y\sqcup L$ with
$L=\{x_0,\ldots,x_{k-1}\}$ and $k\geq2$. Suppose $(A,S)$ is a
parabolic marking and $\sigma\in\Aut(A)$ fixes every $y\in Y$ and
sends $x_j$ to $x_{j+1}$, with subscripts modulo $k$.
Then $\sigma^k=\id$ since it fixes every generator. Set
\[
 G=A\rtimes C_k=\langle A,t\mid t^k=1,\ t^{-1}at=\sigma(a)
 \ (a\in A)\rangle,
 \qquad T=Y\sqcup\{x_0,t\}.
\]
This notation fixes the conjugation convention. The quotient onto $C_k$
proves that $t$ has exact order $k$. The set $T$ generates $G$, since
$x_j=t^{-j}x_0t^j$.

For $R\subseteq T$ define a subset $E(R)$ of $S$ by
\[
 E(R)=\begin{cases}
 R,&t\notin R,\\
 R\cap Y,&t\in R,\ x_0\notin R,\\
 (R\cap Y)\cup L,&t,x_0\in R.
 \end{cases}
\]
In the first line $R$ is read inside $Y\sqcup\{x_0\}\subseteq S$.
\begin{theorem}\label{C-thm:cyclic}
For every $R\subseteq T$ the subgroup generated by $R$ is
\begin{equation}\label{C-eq:all-subgroups}
 G(R)=\begin{cases}
 A(E(R)),&t\notin R,\\
 A(E(R))\rtimes C_k,&t\in R.
 \end{cases}
\end{equation}
The marking $(G,T)$ has the intersection and coset Helly properties.
Consequently its standard coset system is flag-transitive, residually
connected and firm, and its coset complex is flag. The product identity
of Lemma~\ref{C-lem:product} holds for every three generated parabolics.
Its panels have size $k$ at $t$ and $|\langle s\rangle|$ at every
$s\in Y\sqcup\{x_0\}$.
\end{theorem}
\begin{proof}
If $t\notin R$, the assertion is the embedding of $A$ in $G$.
If $t\in R$ but $x_0\notin R$, it fixes every named generator in
$R\cap Y$, giving the second case. If both belong to $R$, conjugating
$x_0$ produces exactly the whole orbit $L$, and this subgroup of $A$
is invariant under $t$. This proves (\ref{C-eq:all-subgroups}) without
an order calculation or finiteness assumption on $A$.

If both subsets contain $t$, the two $A$-parts intersect by
(\ref{C-eq:ip}); the orbit $L$ is present in their intersection precisely
when $x_0$ belongs to both subsets. Thus their intersection is the
subgroup for $R\cap Q$. If neither contains $t$, use (\ref{C-eq:ip})
inside $A$. In the mixed case an element of the intersection has trivial
$C_k$ coordinate. The smaller support contains at most $x_0$ from $L$,
so intersecting its support with the other support again gives exactly
$E(R\cap Q)$. This proves full generated-subgroup intersection.

It remains to prove Helly, which does not follow from subgroup
intersection alone. The maximal parabolics just computed are
\begin{equation}\label{C-eq:maximal}
 P_t=A(Y\cup\{x_0\}),\quad
 P_{x_0}=A(Y)\rtimes C_k,\quad
 P_y=A(S\setminus\{y\})\rtimes C_k\ (y\in Y).
\end{equation}
The $A$-supports in the latter groups are invariant under $C_k$.
Take pairwise-intersecting left cosets of these groups. Repeated types
coincide. If a $P_t$ coset occurs, translate its $C_k$ coordinate to
$1$. All other cosets project onto $C_k$ and have representatives
$(a,1)$. Their fibers at coordinate $1$ are cosets of the indicated
standard parabolics in $A$. Two cosets of the latter kind intersect
if and only if their $A$ fibers intersect, because they are Cartesian
products with $C_k$; incidence with the $P_t$ coset likewise takes
place in this fiber. Thus the $A$ fibers meet pairwise and base Helly
gives a common representative. If no $P_t$ coset occurs, the same
argument applies at coordinate $1$ without translation. This proves
Helly for maximal parabolics. Replace a coset of any generated
parabolic by its maximal-parabolic cosets, using the intersection
property just proved, to get Helly for all of them. Finally apply
Lemma~\ref{C-lem:product}. Every named generator is nontrivial and the
panel sizes are at least two, proving firmness.
\end{proof}
This is the rank-one cyclic case of the orbit construction used in
Piedade--Tranchida's papers on twisting. Its proof is included in full
so that its geometric and presentation conclusions have a complete local proof. It does not assert that an arbitrary Shephard marking has
these properties.

\subsection{The involutory presentation and the doubled braid}
Write $\W_m(u,v)$ for the alternating word of \emph{length} $m$,
starting with $u$. It is not the ordinary power $(uv)^m$.
Let $A$ be the Artin group on $Y\sqcup\{x,x'\}$, with Artin
relations $\mathcal R_Y$ on $Y$, identical labels $m_y$ on
$(x,y)$ and $(x',y)$, and label $m$ on $(x,x')$.
All labels lie in $\{2,3,\ldots\}\cup\{\infty\}$; an infinite
label means no relation. Interchanging $x,x'$ is an automorphism of
the presented group, fixes $Y$, and has order two because the names
are distinct by the retained Artin parabolic theorem.

\begin{lemma}\label{C-lem:braid}
If $t^2=1$, then for every $m\geq1$,
\[
 \W_m(x,txt)=(xt)^m t^{m\bmod2},\qquad
 \W_m(txt,x)=(tx)^m t^{m\bmod2}.
\]
Consequently the braid of length $m$ on $x,txt$ is equivalent to the
braid of length $2m$ on $x,t$.
\end{lemma}
\begin{proof}
For $m=2r$, the two words are $(xtxt)^r=(xt)^{2r}$ and
$(txtx)^r=(tx)^{2r}$. For $m=2r+1$, append respectively $x$ and
$txt$ to those words, obtaining $(xt)^{2r+1}t$ and
$(tx)^{2r+1}t$. Cancellation of the common final $t$ in the odd case
proves the equivalence. The alternating words of length $2m$ are
exactly $(xt)^m$ and $(tx)^m$.
\end{proof}

\begin{theorem}\label{C-thm:fold}
The semidirect product which interchanges $x,x'$ has presentation
\begin{align*}
 \langle Y,x,t\mid {}&\mathcal R_Y,\ t^2=1,\ [t,y]=1\ (y\in Y),\\
 &\W_{m_y}(x,y)=\W_{m_y}(y,x)\ (m_y<\infty),\\
 &\W_{2m}(x,t)=\W_{2m}(t,x)\ (m<\infty)\rangle.
\end{align*}
This is a Shephard presentation, and its named marking satisfies every
conclusion of Theorem~\ref{C-thm:cyclic}.
\end{theorem}
\begin{proof}
Map the displayed generators to the semidirect product. They satisfy
every relator: the last by Lemma~\ref{C-lem:braid}; the others follow
from the defining action and the original Artin relations. They generate,
since $x'=txt$. Conversely, in the displayed presentation set
$x'=txt$. Conjugating a mixed braid by $t$, which fixes $y$, gives
the braid on $x',y$. Lemma~\ref{C-lem:braid} supplies the braid on
$x,x'$. The relations on $Y$ are unchanged. Thus the original Artin
presentation maps into this group, and conjugation by $t$ implements
its transposition of the two leaves. The semidirect-product universal
property gives the inverse map. Both composites fix all named
 generators, so these maps are inverse. This establishes the exact
 group and marking, not merely a quotient action. The cyclic theorem
with $k=2$ now applies. An edge absent from the drawing has label two,
so all commutations $[t,y]=1$ remain part of the presentation.
\end{proof}

\subsection{Cyclic amalgamation}
Let $A_0$ be an Artin group on $Y\sqcup\{x\}$ and
$C=A_0(Y)$. The parabolic embedding theorem identifies $C$ with the
Artin group on the induced graph on $Y$. Form the amalgam $M$ of
$k\geq2$ copies of $A_0$ along this common $C$. Its presentation has
$Y$ and $x_0,\ldots,x_{k-1}$, the relations $\mathcal R_Y$, and the
same braid between $x_j$ and each $y\in Y$. There are no relations
between distinct $x_j$. This follows directly from the universal
property of the amalgam; the embeddings of each factor are the
classical normal-form theorem. The resulting presentation is itself
Artin, with all distinct leaf-pair labels infinite.

\begin{theorem}\label{C-thm:amalgam}
Cyclic permutation of the copies defines an automorphism $\sigma$ of
$M$. With $t^{-1}x_jt=x_{j+1}$, its semidirect product has presentation
\[
 M\rtimes C_k\cong
 \langle Y,x_0,t\mid\mathcal R_Y,\ t^k=1,\ [t,y]=1\ (y\in Y),\
 \W_{m_y}(x_0,y)=\W_{m_y}(y,x_0)\ (m_y<\infty)\rangle.
\]
There is no relation on the pair $x_0,t$. This Shephard marking has
full generated-parabolic intersection and the product identity,
and its standard coset system is flag-transitive, residually connected
and firm with flag coset complex. Its panels have sizes $k$ and
infinity.
\end{theorem}
\begin{proof}
The permutation preserves the Artin relators and has order $k$ on
the distinct leaf names. The displayed generators map to the
semidirect product and generate it. In the converse direction define
$x_j=t^{-j}x_0t^j$. Since $t$ fixes $Y$, conjugation by $t^j$
transports every mixed braid to the required braid on $x_j,y$.
The relation $t^k=1$ closes the cycle. The Artin presentation of
$M$ and then its semidirect product therefore map into the displayed
group. Both composites fix $Y,x_0,t$, proving inverse maps. Apply
Theorem~\ref{C-thm:cyclic} to the Artin marking of $M$.
\end{proof}

\section{Marked permutation modules and changes of coordinates}
Let $r\geq2$, let $V$ be an elementary abelian $2$-group, regarded as an
$\mathbb F_2[S_r]$-module, and let $z\ne0$ be fixed by the stabilizer of
$1$ in $S_r$. Write $s_i=(i\ i+1)$ for $1\leq i<r$. For each $j$ put
$z_j=\pi z$ for any $\pi$ taking $1$ to $j$. This is well-defined:
two such permutations differ on the right by an element fixing $1$.
Let $V_z$ be the span of these $r$ vectors. In $V_z\rtimes S_r$ put
\[
 \rho_0=(z,1),\qquad \rho_i=(0,s_i)\quad(1\leq i<r).
\]
The multiplication convention is $(v,\pi)(w,\sigma)=(v+\pi w,\pi\sigma)$.

\begin{theorem}[A marked module]\label{M-thm:module-marking}
The displayed marks form a string C-group of rank $r$, with generated
group $V_z\rtimes S_r$. This assertion requires no linear independence
of the orbit vectors $z_1,\ldots,z_r$.
\end{theorem}
\begin{proof}
All marks are nonidentity involutions. The marks $\rho_i$ with $i\geq1$
generate the indicated copy of $S_r$; their conjugates of $\rho_0$
generate $V_z$. If $i\geq2$, then $s_i$ fixes $1$ and hence $z$, so
$\rho_0$ commutes with $\rho_i$. The other string commutations are
those of adjacent transpositions.

We prove every subgroup intersection explicitly. For a set of mark
indices $D\subseteq\{0,\ldots,r-1\}$ let $E=D\setminus\{0\}$ and
$P_E=\langle s_i:i\in E\rangle$. This is the product of the symmetric
groups on the connected components of the path with retained edges $E$:
path transpositions generate each component and move no point outside
it. Its orbit on $1$ is an initial segment $[k(E)]$, where $k(E)$ is
one plus the length of the initial run $1,2,\ldots$ contained in $E$.
Set
\[
 W_D=\begin{cases}
 \langle z_1,\ldots,z_{k(E)}\rangle_{\mathbb F_2},&0\in D,\\
 0,&0\notin D.
 \end{cases}
\]
The group generated by $D$ is exactly $W_D\rtimes P_E$. Indeed, all
its vector generators are conjugates of $z$ under $P_E$, and those
conjugates are precisely the vectors in the displayed span. The
other path components fix $z$ and do not introduce any further vector.

For two edge sets $E,F$, a permutation preserving each block of both
component partitions preserves each of their pairwise intersections. Those
intersections are exactly the components of the path retaining
$E\cap F$: two vertices are in the same component precisely when every
edge of their unique interval belongs to both sets. Thus
$P_E\cap P_F=P_{E\cap F}$. Also
$k(E\cap F)=\min\{k(E),k(F)\}$. The vector spaces spanned by initial
segments form a nested chain, even when some orbit vectors are dependent.
Consequently $W_D\cap W_{D'}=W_{D\cap D'}$. Intersecting inside the
Cartesian set $V_z\times S_r$ proves the full C-group property.
\end{proof}

The elementary coordinate change below applies to any group action on
an abelian group, not only to permutation modules.
\begin{lemma}[Cocycle change of coordinates]\label{M-lem:shear}
Let $Q$ act on an abelian group $V$, written additively. Suppose
$c:Q\to V$ satisfies
\[
 c(gh)=c(g)+g c(h).
\]
Then $\Phi_c(v,g)=(v+c(g),g)$ is an automorphism of $V\rtimes Q$,
with inverse $\Phi_{-c}$. It fixes $V$ pointwise and carries any marked
subgroup intersection identity to the corresponding identity on the
image marks. In characteristic two it is an involution.
\end{lemma}
\begin{proof}
The cocycle identity with $g=h=1$ gives $c(1)=0$. For two elements,
\[
 \Phi_c((v,g)(w,h))=(v+gw+c(g)+gc(h),gh)
 =\Phi_c(v,g)\Phi_c(w,h).
\]
Addition of $-c(g)$ in the same coordinate is an inverse. The last
assertions follow from $c(1)=0$, from bijectivity, and from
$c(g)+c(g)=0$ in characteristic two. In particular an automorphism
commutes with taking generated subgroups and their intersections.
\end{proof}

\begin{proposition}[Three pairs of signed markings]\label{M-prop:signed}
Take $V=\mathbb F_2^r$ with its natural permutation action. Write
$e_1,\ldots,e_r$ for its basis and $\lambda=e_1+\cdots+e_r$. For any
of the three seeds
\[
 z=\lambda,\qquad z=e_1,\qquad z=\lambda+e_1,
\]
both of the following lists are string C-group markings:
\[
 (z,1),(0,s_1),\ldots,(0,s_{r-1}),
 \qquad
 (z,1),(\lambda+e_1+e_2,s_1),\ldots,
       (\lambda+e_{r-1}+e_r,s_{r-1}).                 \tag{*}
\]
The common abstract groups for these three pairs are respectively
$C_2\times S_r$, $C_2\wr S_r$, and
\[
 \begin{cases}
 C_2\wr S_r,&r\text{ even},\\
 V^{\mathrm{even}}\rtimes S_r,&r\text{ odd},
 \end{cases}
 \qquad V^{\mathrm{even}}=\{v:\sum_j v_j=0\}.
\]
In the last line $r\geq3$ when odd. The statement includes the odd-rank
second list for the third seed.
\end{proposition}
\begin{proof}
Each seed is nonzero and fixed by the stabilizer of $1$, so the first
list is covered by Theorem~\ref{M-thm:module-marking}. To identify its vector
span, only the third seed needs explanation. Its orbit vectors are
$\lambda+e_j$. A relation with coefficients $b_j$ says
$b=(\sum_j b_j)\lambda$. For even $r$ this forces $b=0$; for odd $r$
the only nonzero relation is $b=\lambda$. In the odd case every orbit
vector has even coordinate sum, so their $(r-1)$-dimensional span is
exactly $V^{\mathrm{even}}$.

Let $a\in V$ have alternating coordinates, so that $a_i+a_{i+1}=1$.
Let $\epsilon:S_r\to\mathbb F_2$ be the sign homomorphism. Define
\[
 c(\pi)=\epsilon(\pi)\lambda+a+\pi a.
\]
The first summand is a cocycle because $\lambda$ is fixed and
$\epsilon$ is additive. The other two are a coboundary:
$a+\pi\sigma a=(a+\pi a)+\pi(a+\sigma a)$.
For $s_i$ the result is $c(s_i)=\lambda+e_i+e_{i+1}$.
Lemma~\ref{M-lem:shear}, in the full ambient $V\rtimes S_r$, therefore
carries the first list in (*) to the second, and fixes its first mark.
It need not preserve the first list's named subgroup as a subset; it
is an isomorphism onto the subgroup generated by the second list.
This distinction matters for odd $r$, where the image copy of $S_r$
may have odd vector coordinates. It does not affect the intersection
property or the asserted abstract group.
\end{proof}

\paragraph{Source correspondence and attribution.}
The ordinary lists in Proposition~\ref{M-prop:signed}, in the order of
its three seeds, are graph (1) of Table 9 and graphs (13), (14) of the
cited Fernandes--Piedade catalogue. The changed-coordinate lists are
graph (2) of Table 9 and graphs (15), (16), respectively. In the bound
arXiv:2407.16003v2 edition the latter four diagrams are in Table 6,
physical page 25; the target paper variously calls that table 5 or 6.
The catalogue's Proposition 4.7 and Table 6 concern $r=n/2\geq7$;
the algebraic marking theorem proved here holds for every $r\geq2$.
Here a vector $e_j$ flips the two points in block $j$, and $s_i$
interchanges adjacent blocks without flips. Thus the vector
$\lambda+e_i+e_{i+1}$ means precisely the vertical $i$-edges on every
block except $i$ and $i+1$, as shown in those diagrams.

The known family statements are Piedade--Tranchida~\cite[Proposition 5.1]{N006}.
The odd-rank graph (16) is expressly left unverified by their method
in the paragraph immediately before that proposition. The formula
above proves that case as well. The source bridge and the full proof were distinctly reviewed in the precursor edition; the present migration retains both. This is not a claim that no previous work settles it. Cocycle
changes of splitting are standard group cohomology; the contribution
here is their explicit use to identify the paired markings, together
with the nested-orbit-span proof which tolerates the odd-rank relation.
No source-specific theorem is presumed in either proof.

\section{Attribution and retained interfaces}
The target constructions are due to Claudio Alexandre Piedade and Philippe Tranchida~\cite{N005,N006,N007}. The tree support and local-coset arguments, indexed qualifications, cyclic presentation maps and module-coordinate proofs here reorganize the stated results and supply explicit corrections; no historical priority is asserted. Normal forms and the Bass--Serre tree construction are retained from classical combinatorial group theory, including Lyndon--Schupp, \emph{Combinatorial Group Theory}, Chapters I and IV, and Serre, \emph{Trees}. The Artin and Coxeter parabolic interfaces are stated with exact locators in the cyclic section. Their complete external proofs are not reproduced and their costs are not set to zero.

The existing \texttt{OrbitTwisting}, \texttt{PrefixSpans} and \texttt{CocycleShear} Lean files certify only their declared algebraic interfaces. They do not formalize the new tree, HNN, cyclic Helly, reflection or affine arguments. This edition adds no Lean declaration. Mathematical and formal verification are separate boundaries.
\begin{thebibliography}{9}\raggedright
\bibitem{N005} C. A. Piedade and P. Tranchida, \emph{From Group Operations to Geometric Structures: Amalgamations, HNN-Extensions, and Twisting in Coset Geometries}, arXiv:2505.24662v1 (2025).
\bibitem{N006} C. A. Piedade and P. Tranchida, \emph{The geometry of wreath and semi-direct products}, arXiv:2605.04646v2, revised edition (2026).
\bibitem{N007} C. A. Piedade and P. Tranchida, \emph{Coset geometries acting on their elements: The $j$-diagonals twisting}, arXiv:2608.23223v1 (2026).
\bibitem{FP} M. E. Fernandes and C. A. Piedade, \emph{Regular polytopes of rank $n/2$ for transitive groups of degree $n$}, arXiv:2407.16003v2 (2025); J. Algebra 689 (2026), 256--296. The existing signed-family adapters use Tables 6 and 9 of the bound arXiv edition.
\end{thebibliography}
\end{document}
